% Guia do professor — Capítulo 15: Perigos de Tempestades
\section{Capítulo 15 --- Perigos de Tempestades}

\subsection*{Visão geral e ritmo}

O capítulo mais ``vendável'' do curso — e uma coleção de física
clássica no extremo: balanço peso--arrasto (granizo), termodinâmica da
evaporação (downburst), eletrostática (raios) e ciclostrofia
(tornado). Cada perigo fecha com números de segurança reais.
\textbf{Ritmo: 2 aulas de 2\,h.}

\begin{itemize}
  \item \textbf{Aula 1} — granizo (balanço no quadro, biografia da
        pedra, derretimento na queda) e downbursts (DCAPE, seco ×
        úmido, microburst e aviação). Casos~1--2.
  \item \textbf{Aula 2} — relâmpagos (o capacitor no quadro,
        flash-to-bang, 30--30) e tornados (Rankine, queda de pressão,
        escala EF, lado perigoso no HS). Casos~3--4.
\end{itemize}

\subsection*{Objetivos de aprendizagem}

\begin{enumerate}
  \item calcular o updraft mínimo para cada tamanho de granizo;
  \item definir DCAPE e prever a intensidade de downbursts a partir da
        umidade sub-nuvem;
  \item descrever a eletrificação por colisões e usar as regras de
        segurança (3\,s/km, 30--30);
  \item integrar o equilíbrio ciclostrófico do Rankine e obter
        $\Delta P = \rho v_{max}^2$;
  \item relacionar a escala EF com a pressão dinâmica
        $\tfrac12\rho v^2$.
\end{enumerate}

\subsection*{Pré-requisitos a reativar}

Updraft e CAPE (Cap.~14); riming e fase mista (Cap.~7); $\theta_w$ e
evaporação (Cap.~4); equilíbrio ciclostrófico (Cap.~10); esticamento
de vorticidade (Cap.~13); refletividade e dual-pol (Cap.~8).

\subsection*{Armadilhas conceituais frequentes}

\begin{itemize}
  \item \textbf{``Granizo cai porque ficou pesado''}: cai quando
        $w_T(D)$ VENCE o updraft — a mesma pedra estaria voando numa
        supercélula mais forte.
  \item \textbf{Downburst exige chuva forte}: o seco cai de virga com
        eco fraco — o mais perigoso é o que não se vê.
  \item \textbf{Raio ``procura o mais alto''}: o líder desce por
        passos de campo máximo; altura ajuda mas não decide sozinha —
        daí a regra do abrigo, não da estatística.
  \item \textbf{Tornado como ``mini-ciclone''}: sem Coriolis! É
        ciclostrófico e pode girar para qualquer lado; o giro
        predominante vem do mesociclone-mãe.
  \item \textbf{Lado perigoso}: no HS é o ESQUERDO da trajetória
        (giro horário $+$ translação) — espelho dos livros do HN.
\end{itemize}

\subsection*{Demonstração de baixo custo}

Gerador eletrostático (Wimshurst/van de Graaff) para faíscas e o
conceito de ruptura; vórtice de garrafa PET (dois litros + água) para
o Rankine com ``olho'' visível. Vídeos: microburst de Andrews AFB
(anemômetro de 67\,m/s) e granizo em câmera lenta.

\subsection*{Problemas sugeridos do livro (exercícios do Cap.~15)}

Aquecimento: $w_T$ para tamanhos dados; flash-to-bang de tempos
medidos. Aprofundamento: DCAPE de sondagens; $\Delta P$ de tornados
com $v_{max}$ dados. Síntese: ``por que o corredor sul do Brasil /
norte da Argentina tem os MCCs mais elétricos do planeta?''.
\emph{Confira a numeração na sua cópia (v1.02b).}

\subsection*{Uso do notebook \texttt{cap15\_perigos.ipynb}}

\begin{center}
\begin{tabular}{lll}
\hline
Caso & Conteúdo & Momento sugerido \\
\hline
1 & $w_T(D)$ do granizo com marcos (ervilha$\to$laranja) & Aula 1 \\
2 & downburst: queda com déficit evaporativo & Aula 1 \\
3 & capacitor da tempestade: dente de serra de raios & Aula 2 \\
4 & Rankine: $v(r)$ e $\Delta P$ por EF & Aula 2 \\
\hline
\end{tabular}
\end{center}

\subsection*{Exercícios complementares e soluções}

Nas notas: T1--T5 e N1--N4. Soluções em \texttt{cap15\_solucoes.ipynb}
(\emph{não distribuir}): T1 e T4 com \texttt{sympy}; N1 e N4 são os
mais ricos (derretimento na queda; lado perigoso no HS). Pares:
T1$\leftrightarrow$N1, T2$\leftrightarrow$N2, T4$\leftrightarrow$N4.

\subsection*{Conexões com o resto do curso}

Fábrica (Cap.~14); microfísica do gelo (Cap.~7); radar e dual-pol
(Cap.~8); ciclostrofia (Cap.~10); esticamento (Cap.~13); ventos fortes
não-convectivos para contraste (Cap.~17); climatologia de severos num
clima mais quente (Cap.~21).
